%%%%%%%%%%%%%%%%%%%%%%%%%%%%%%%%%%%%%%%%%
% MUW Poster
% LaTeX Template
% Version 1.0 (31/08/2016)
% (Based on Version 1.0 (31/08/2015) of the Jacobs Portrait Poster
%
% License:
% CC BY-NC-SA 3.0 (http://creativecommons.org/licenses/by-nc-sa/3.0/)
%
% Created by:
% Nicolas Ballarini, CeMSIIS, Medical University of Vienna
% nicoballarini@gmail.com
% http://statistics.msi.meduniwien.ac.at/
%%%%%%%%%%%%%%%%%%%%%%%%%%%%%%%%%%%%%%%%%


\def\footer#1{\def\insertfooter{#1}}
%--------------------------------------------------------------------------------------
%	PACKAGES AND OTHER DOCUMENT CONFIGURATIONS
%--------------------------------------------------------------------------------------

\documentclass[final]{beamer}



\usepackage[scale=1.150]{beamerposter} % Use the beamerposter package
\usetheme{MUWposter} % Use the MUWposter theme supplied with this template

% Include a logo of your project if desired
\logo{\pgfputat{\pgfxy(-12,108)}{\pgfbox[center,base]{\includegraphics[width=8cm]{pbn.png}}}}  

\usepackage{multicol}
\usepackage{array}
%The following two are column definitions for the aknowledgements section
\newcolumntype{L}{>{\arraybackslash}m{22cm}}
\newcolumntype{S}{>{\arraybackslash}m{5cm}}
\usepackage{pgf}  
\usepackage{mathtools}
\usepackage{amsmath, amsthm, amssymb, amsfonts}
\usepackage{exscale}
\usepackage{xcolor}
\usepackage{ushort}
\usepackage{setspace}
\usepackage[square,numbers]{natbib}
\usepackage{url}
\bibliographystyle{abbrvnat}
\renewcommand{\vec}[1]{\ushort{#1}}
\renewcommand{\vec}[1]{\mathbf{#1}}
\definecolor{greenMUW}{RGB}{60,191,174}
\definecolor{blueMUW}{RGB}{17,29,79}
\definecolor{skinMUW}{RGB}{254,228,217}
\definecolor{hellblauMUW}{RGB}{95,180,229}

%-----------------------------------------------
%  START Set the colors
%  Uncomment to apply colors you want to use.
%-----------------------------------------------
\colorlet{themecolor}{greenMUW}
\usebackgroundtemplate{\includegraphics{MUW_green.pdf}}

%\colorlet{themecolor}{skinMUW}
%\colorlet{themecolor}{blueMUW}
%\usebackgroundtemplate{\includegraphics{MUW_skin.pdf}}

%%\colorlet{themecolor}{blueMUW}
%\colorlet{themecolor}{hellblauMUW}
%\usebackgroundtemplate{\includegraphics{MUW_hellblau.pdf}}
%-----------------------------------------------
%  END Set the colors
%-----------------------------------------------


%-----------------------------------------------
%  START Set the width of the columns
%-----------------------------------------------
\setlength{\paperwidth}{33.1in} % A0 width: 46.8in
\setlength{\paperheight}{46.8in} % A0 height: 33.1in
\newlength{\sepmargin}
\newlength{\sepwid}
\newlength{\onecolwid}
\newlength{\twocolwid}
\newlength{\threecolwid}

% The following measures are used for 2 columns
\setlength{\sepmargin}{0.055\paperwidth} % Separation width (white space) between columns
\setlength{\sepwid}{0.03\paperwidth} % Separation width (white space) between columns
\setlength{\onecolwid}{0.43\paperwidth} % Width of one column
\setlength{\twocolwid}{0.9\paperwidth} % Width of two columns

%-----------------------------------------------------------
% The following measures are used for 3 columns
%\setlength{\sepmargin}{0.06\paperwidth} % Separation width (white space) between columns
%\setlength{\sepwid}{0.02\paperwidth} % Separation width (white space) between columns
%\setlength{\onecolwid}{0.28\paperwidth} % Width of one column
%\setlength{\twocolwid}{0.58\paperwidth} % Width of two columns
%\setlength{\threecolwid}{0.88\paperwidth} % Width of three columns
%\setlength{\columnsep}{30pt}

%-----------------------------------------------
%  END Set the width of the columns
%-----------------------------------------------


%--------------------------------------------------------------------------------------
%	TITLE SECTION 
%--------------------------------------------------------------------------------------
\setbeamertemplate{title}[left]
\setbeamertemplate{frametitle}[default][left]
%\setmainfont{Georgia}

\title{Creating Knowledge graphs from the literature: the case of health resilience in Green Building Neighbourhoods} % Poster title

\author{Luc Jonveaux, Filip Kucera, Jorge, Marta, Kathrine, Sissa} % Author(s)

\institute{The Integrator-centric approach for realising innovative energy efficient buildings in connected sustainable green neighbourhoods project} % Institution(s)
%--------------------------------------------------------------------------------------



\begin{document}
%% https://cordis.europa.eu/project/id/101037075/
  \addtobeamertemplate{block end}{}{\vspace*{1ex}} % White space under blocks
  \addtobeamertemplate{block alerted end}{}{\vspace*{0ex}} % White space under highlighted (alert) blocks
  \setlength{\belowcaptionskip}{2ex} % White space under figures
  \setlength\belowdisplayshortskip{1ex} % White space under equations
  
  
  \begin{frame}[t] % The whole poster is enclosed in one beamer frame

      \begin{columns}[t] % The whole poster consists of two major columns
	  
      \begin{column}{\sepmargin}\end{column}
      
	    \begin{column}{\onecolwid} % The first column


		  \begin{block}{Context}
          %\begin{multicols}{2}
The Horizon 2020\href{https://doi.org/10.3030/101037075}{PROBONO project} (Grant agreement ID: 101037075) aims at demonstrating "\textit{strong examples of how Green Building Neighbourhoods (GBNs) technological and social innovations can be applied, with a vision focused on building infrastructure and a renewed focus on people and sustainability, taking full advantage of digitization and smart technologies for the benefit of society}". Task 3.5 of this projects  aims at reviewing "Interventions to mitigate diseases outbreaks".

In this document, we summarize how we used new technologies to consolidate a knowledge graph of this topic, based on the literature, to demonstrate the feasibility of building a body of knowledge pertaining to a certain domain, with a specific angle.
\end{block}

\begin{block}{Disclaimer}
Please note that this document is a research-based exploration compiled by knowledge management researchers, not medical professionals. Our findings are presented with the intention to inform and contribute to the dialogue on public health strategies. They are indicative and should serve as a preliminary guide. We encourage all readers to consult with qualified health professionals for expert advice and to confirm and enrich these insights.

          %\end{multicols}
          \end{block}
          
          \begin{block}{Objective}
The main objective of the task is to review the scientific literature to identify key risks, stakeholders, technologies and mitigations measures both at building and neighbourhood scales.

          \end{block}
          
          
          \begin{block}{Results}
          %\begin{multicols}{2}

We started with defining a basic ontology based on classes of interest for mitigation measures, namely 'Risks', 'Mitigations', 'People', and 'Technologies'. These constitute an initial body of knowledge, which is then used to build 'Blueprints', possible interventions to mitigate diseases outbreaks.

To date, the Knowledge Graph (v0.3) contains information on 376 articles, from which were programmaticaly derived 3418 risks, 5295 mitigation measures, 2640 stakeholders and 3915 technologies. The team used these to build 24 blueprints manually, and automated the production of 50 others.

We hope that releasing the knowledge graph under an open-source license will drive use of this knowledge graph and that health professionals can use this to derive useful, professionally-approved mitigation measures.

          %\end{multicols}
          \end{block}
         
         \end{column}
                  
                  
        % SEparator column
         \begin{column}{\sepwid}  \end{column}
         
         
         
         
         \begin{column}{\onecolwid} %The second column
         
         \begin{block}{Technical details}
          %\begin{multicols}{2}
          The present knowledge graph has been created using new tools, helping to streamline, faster and more consistently, information from the litterature:
           \begin{itemize}
\item Parsing of the litterature was done with GROBID. This provided structured text (XML) from the article PDFs. 
\item The data was processed and structured in an RDF that was produced with the pyowlready2 library.
\item For vector embedding, we used Pinecone because of the integration with Langchain and ease of use.

 \end{itemize}
Extracting information was done using different solutions:
           \begin{itemize}
\item Using NLTK to extract topics

\item Spacy for entity recognition, and corefeeringing


\item Using OpenAI API (both using GPT3.5 and GPT4 endpoints). This was done using the request library or the langchain library.

\item Ex.

\item X.
          %\end{multicols}
                     \end{itemize}
                     We still hope that
                                \begin{itemize}
\item We plan on using Neo4j or similar to continually review and enrich the knowledge graph.
     %\end{multicols}
                     \end{itemize}
          \end{block}
          
      \end{column}
      
      \begin{column}{\sepmargin} \end{column}
      \end{columns} 
       
      \begin{columns}[t] % Split up the two columns wide column again
      
      \begin{column}{\sepmargin} \end{column}
        \begin{column}{\onecolwid} % The first column
			\begin{block}{\large Acknowledgements}
                    Thank you to the PROBONO team.
                    %\begin{center}
				%		\begin{tabular}{SL}
			%				\includegraphics[width=\linewidth]{Flag_of_Europe.png}  &
			%				\footnotesize This project has received funding from the  grant agreement No 11111.
			%			\end{tabular}
			%		\end{center}
				\end{block}	
                \vspace*{-0.9cm}
				\begin{alertblock}{\large Contact Information}
                \vspace*{-0.5cm}
					\begin{footnotesize}
					\begin{itemize}
						\item \href{mailto:luc.jonveaux@mottmac.com}{luc.jonveaux@mottmac.com}
						\item \href{http://www.mottmac.com/}{www.mottmac.com}
					\end{itemize}
					\end{footnotesize}	
					
				\end{alertblock}
		    \end{column} % End of the first column
			\begin{column}{\sepwid}\end{column} % Empty spacer column
			\begin{column}{\onecolwid} % Begin a column 
              \begin{block}{\large References}
			  %\vspace*{-0.5cm}
              	\nocite{*} % Insert publications even if they are not cited in the poster
					
                    	%\bibliographystyle{plainurl}
						n/a
      
				\end{block} 
    \begin{block}{\large Note(s)}
			  %\vspace*{-0.5cm}
              	\nocite{*} % Insert publications even if they are not cited in the poster
					%{\footnotesize
                    	%\bibliographystyle{plainurl}
						This PDF can be renamed as a ZIP file to extract source code for this document and for the app.
      % @LJ: For this, see https://github.com/ansemjo/truepolyglot
      %}
				\end{block} 
			\end{column} % End of the second column
            
			\begin{column}{\sepmargin}\end{column} % Empty spacer column
            
\end{columns} % End of all the columns in the poster


\end{frame} % End of the enclosing frame
	
\end{document}
